%% 3_method
\section{The Geometry a Convolution Cannot Form}
\label{sec:method}

\begin{figure*}[t]
\centering
\includegraphics[width=\textwidth]{fig/concept.pdf}
\caption{\textbf{The missing input.} Given a building layout and a transmitter
(left, red star), the occlusion field of \cref{eq:occ} integrates building occupancy
along each transmitter-to-pixel ray and recovers the radial shadow structure that
shapes the true field (right). Replacing the transmitter anchor with a minimum over
measurement locations (\cref{eq:occmeas}, third panel) collapses the field toward
zero, with a scene mean of $0.084$ against $0.298$, and destroys exactly the
directional structure that makes the channel useful.}
\label{fig:concept}
\end{figure*}

\subsection{Reconstruction from walks, not from samples}

A radio field is a dense image that no sensor can capture in one exposure. We write
$R : \Omega \to \mathbb{R}$ for the field over a discrete spatial domain
$\Omega \subset \mathbb{Z}^2$ of $256 \times 256$ cells, with values in dBm. A
reconstructor observes $R$ only at a small measured subset
$\mathcal{M} \subset \Omega$ and must predict it everywhere, including in the free
space it never visited.

The acquisition process decides what $\mathcal{M}$ looks like, and that is where
this problem departs from the sparse-to-dense tasks vision usually studies.
Crowdsensed measurements arrive along walks. We therefore draw $k=3$ trajectories of
$100$ points each per scene, giving $300$ observations, which covers a small
fraction of the free space and leaves whole regions unobserved rather than
uniformly thinned. Accuracy is reported on free-space cells that were never
measured, since those are the cells a deployment actually needs and the ones
interpolation cannot reach.

\subsection{A line integral the backbone cannot build}

Signal strength at a point is governed less by distance than by what the signal
passed through on the way. For a transmitter at $\tau \in \Omega$ and a target cell
$p \in \Omega$, define the occlusion field as the fraction of the segment
$\overline{\tau p}$ that lies inside building material,
%
\begin{equation}
O(p) \;=\; \frac{1}{N}\sum_{i=1}^{N} B\!\left(\tau + \tfrac{i}{N}\,(p - \tau)\right),
\qquad N = 32,
\label{eq:occ}
\end{equation}
%
where $B : \Omega \to \{0,1\}$ is the building raster and intermediate positions are
resolved by bilinear sampling, which makes $O$ differentiable in $\tau$. The
quantity is a soft shadow depth. It is zero along clear line of sight, rises as the
ray clips a corner, and saturates when the path is buried in material.

\Cref{eq:occ} is a line integral between two points that may be two hundred cells
apart, and that is exactly the operation a convolutional stack is poorly equipped to
perform. A convolution aggregates over a compact neighbourhood, so forming $O$
internally would require the receptive field to span the scene and then to learn a
path-dependent selection within it. The field is fully determined by inputs the
network already receives, namely the building raster and the transmitter position,
yet it is not a function any local operator computes cheaply. We therefore compute
it once in closed form and append it as a single input channel. The cost is
$432$ parameters in the first convolution of the single network, or $576$ spread
over the two stages of the cascade, and no change anywhere else.

\subsection{Two backbones, one variable}

We study two backbones so that the effect of the channel can be separated from the
effect of capacity. The first is a single encoder-decoder network with $48$ base
channels and depth $4$. The second is a cascade of two such networks in series, the
second taking the first network's output alongside the original inputs, with the
base width scaled to $0.72$ so the pair roughly matches the single network's
parameter budget. The cascade is the backbone RadioUNet uses, which makes it the
right control for asking whether our gain comes from geometry or from borrowed
architecture.

Everything else is held identical so that the channel is the only thing that moves.
Both backbones receive the sparse measurement raster, its binary mask, a coverage
density map, the building raster and a multi-scale transmitter heatmap, and adding
the occlusion channel is the single variable under study. Dropout is retained at
inference to give a per-pixel predictive spread by Monte Carlo sampling.

\subsection{Dropping the transmitter anchor}
\label{sec:txfree}

\Cref{eq:occ} needs $\tau$, and that is a real liability. Indoor cellular
deployments place boosters that share cell identifiers with the macro network, so
the source of a measured signal is frequently not locatable at all. On our real
measurements only $17.9\%$ of cells yield a transmitter position that passes a
reliability gate, which means a conditioning channel anchored at $\tau$ is
unavailable for roughly four cells in five.

One obvious escape is to anchor the ray somewhere the deployment does know. Let
$\mathcal{A} \subset \mathcal{M}$ be a sample of measurement locations, whose positions
are recorded whether or not the source is identified.
Define the transmitter-free field as the least obstructed path from any measurement
to the target,
%
\begin{equation}
\tilde{O}(p) \;=\; \min_{a \in \mathcal{A}} \; \frac{1}{N}\sum_{i=1}^{N}
 B\!\left(a + \tfrac{i}{N}\,(p - a)\right),
\qquad |\mathcal{A}| = 16 .
\label{eq:occmeas}
\end{equation}
%
The intuition is that a cell visible from some point at which the signal was
observed should behave like a cell in line of sight. \Cref{sec:txfree_result}
reports what happened, which is that the substitution does not work, and
\cref{fig:concept} shows why.
